\documentclass[11pt]{article}
\usepackage[margin=1in]{geometry}
\usepackage[T1]{fontenc}
\usepackage{lmodern,booktabs,longtable,amsmath}
\usepackage[hidelinks]{hyperref}
\usepackage{xurl}
\setlength{\parindent}{0pt}
\setlength{\parskip}{6pt}
\title{CARLA Estimator Smoke Test}
\author{Pan Dynamics Collision Avoidance Research}
\date{October 2, 2026}
\begin{document}
\maketitle

\section{Outcome and scope}
The unchanged production NRMM estimator completed both 400-frame replays of
a newly collected physical CARLA two-vehicle trajectory: continuous target
measurements and a 0.30-second target-measurement dropout. All published
states remained finite, sample times stayed aligned, and all six descriptive
RMSE gates passed after a two-second transient. The existing
\path{tests/onlineNrmmTrackingRuntimeTest.m} suite passed all 31 tests, with
zero failures or incomplete tests (63.4107 seconds of test execution).

This is an estimator functional smoke test using simulator-state measurement
proxies. It is not a native-sensor perception integration test, closed-loop
collision-avoidance test, real-time guarantee, or validation of the theoretical
error certificate. The trajectory violates several default model assumptions;
the relative-position bound is unavailable at 377 of 400 frames in each replay.
No production estimator, controller, gain, or configuration was changed.

\section{Server isolation and capture}
The shared CARLA server on port 2000 was actively used by another session's
multi-route experiment. Read-only inspection was followed by a separate
temporary server on port 2011, using the installed package, null RHI, no sound,
low quality and a nice level of 10. This test issued no world-setting changes,
actor mutations, or ticks to port 2000. The original shared server process
(PID 2454708) was still listening after the temporary server stopped.

The temporary server and Python API both reported CARLA 0.10.0. The map was
\path{Carla/Maps/Town10HD_Opt}. The existing local
\path{simulation/carla_town10_capture.py} supplied vehicle spawning,
waypoint following, sensor synchronization and frame-record helpers. That
legacy recorder requires RGB/LiDAR and lacks the current estimator's GNSS
velocity and target-position fields. A local test adapter omitted rendering
and added these fields without restoring legacy tooling to version control.

The physical ego was \path{vehicle.lincoln.mkz}; the target was
\path{vehicle.dodge.charger}. The existing follower used nominal speeds of
11 and 12 m/s, respectively, a 10 m look-ahead, a 25 m initial waypoint gap,
and spawn index zero. Physics remained enabled. After 200 warm-up frames,
400 consecutive frames were collected at 0.02 seconds per frame. The recorded
sample interval is 0--7.98 seconds; 600 world ticks cover warm-up and recording.
All 400 native IMU/GNSS pairs matched their CARLA frame IDs. The maximum
sample-spacing discrepancy was $4.48\times10^{-10}$ seconds.

After capture, all owned actors were destroyed and the original settings
were restored. The immediate actor lookup still returned cached IDs; a separate
subsequent query confirmed zero vehicles and zero sensors on port 2011.
The owned process (PID 2478952) received SIGTERM only after this verification.
Its launcher then reported a segmentation fault / exit code 139 during
shutdown. Capture and both estimator replays had completed successfully;
port 2011 was confirmed closed. The shutdown fault remains an environment
limitation, not a successful clean process exit. An initial short readiness
probe also timed out during startup; the subsequent capture connection succeeded.

\section{Measurement contract and estimator execution}
The estimator consumed actor-derived planar GNSS position and velocity,
body-frame acceleration, yaw rate and target-center relative position, with
independent bounded noise from NumPy seed 20261002. Two-dimensional noise was
uniform in a disc; yaw-rate noise was uniform in an interval. Bounds matched
the unchanged configuration: 0.12 m position, 0.02 m/s velocity, 0.05 m/s$^2$
acceleration, 0.002 rad/s gyro and 0.12 m target position. An independent
Python audit confirmed every generated error remained within its bound.
Native IMU/GNSS samples were recorded and frame-matched, but were not passed
off as calibrated estimator inputs. No radar detector or YOLO/LiDAR target
extraction was exercised.

Replay reflected inertial and body $y$ coordinates and negated yaw/yaw rate
consistently. Recorded spacing was checked before timestamps were mapped to
the exact integer sample grid. Each call to
\path{onlineNrmmTrackingRuntime('sample',...)} published the state at the
current sample and advanced the runtime by 0.02 seconds. Scoring compared
that publication with same-frame truth, avoiding a one-sample shift.

Initialization used only the first noisy measurement: GNSS position,
velocity/course with the configured single-track sideslip correction, and
target relative position. The initial target was assumed co-moving with zero
acceleration. No ground-truth initialization or future observations were used.
The default gains and domain parameters were retained, including the MnCAV
ego rear-axle distance of 1.714395 m despite the CARLA Lincoln vehicle.
The runtime used four 0.005-second RK4 substeps and the existing compiled
\path{solver/nrmm/nrmmObserverRk4IntervalMex.mexa64}.
The necessary nominal sampled spectral-radius check was 0.9602748642.
The default unbounded ego acceleration-envelope fields serialize as JSON
\texttt{null}; they mean positive infinity in the MATLAB configuration.

\section{Measured results}
Both replays used the same physical capture and noisy measurements.
The dropout run withheld 15 target-position measurements on
$[3.50,3.80)$ seconds, using the runtime's explicit unavailable-detection
contract. All RMSE values below use the 300 samples at $t\geq2$ seconds.
Thresholds were specified in the local replay adapter before execution and
are smoke-test criteria, not accuracy requirements or theoretical bounds.

\begin{center}
\begin{tabular}{lrrr}
\toprule
RMSE & Continuous & Dropout & Smoke limit \\
\midrule
Ego position (m) & 0.01414 & 0.01414 & 0.50 \\
Ego inertial velocity (m/s) & 0.14380 & 0.14380 & 0.50 \\
Ego body velocity (m/s) & 0.24886 & 0.24886 & 0.50 \\
Ego yaw (degrees) & 0.68993 & 0.68993 & 5.00 \\
Target relative position (m) & 0.13541 & 0.13543 & 0.50 \\
Target inertial velocity (m/s) & 1.02063 & 1.01796 & 1.50 \\
\bottomrule
\end{tabular}
\end{center}

Maximum post-transient relative-position errors were 0.26975 m and 0.26960 m.
The dropout run's maximum relative-position error after reacquisition was
0.26960 m. Integration and all six RMSE gates passed for both runs.
Median estimator-call wall times were 2.552 ms and 2.2985 ms; maxima were
429.959 ms and 101.506 ms. These are offline, shared-machine measurements
including cold calls, not a 20 ms deadline guarantee.

\section{Model-domain findings and limits}
The physical route includes sharp turns and a yaw wrap. Ego speed ranged
from 10.3569 to 11.2048 m/s, while target speed ranged from 6.9990 to
11.5185 m/s. Target range was 16.2062--54.7591 m. The following direct
truth audits explain why this trajectory cannot support a certificate claim:
\begin{itemize}
\item Ego yaw rate reached 0.99352 rad/s; 146 frames exceeded the configured
0.30 rad/s maximum.
\item Target speed was below the configured 10 m/s minimum for 95 frames.
\item Target range exceeded the configured 50 m maximum for 12 frames.
\item The ego single-track residual
$\omega_E-v_{Ey}/l_{r,E}$ reached 0.34943 rad/s in magnitude;
269 frames exceeded the configured 0.02 rad/s mismatch allowance.
\end{itemize}
Estimated operating-domain audits were invalid for 254 frames in each run;
relative-position bounds were unavailable for 377 frames. The implementation
continued returning finite estimates as designed. These observations do not
justify widening bounds after the fact or treating the available radii on
other frames as a certification of this CARLA vehicle/route. Native-sensor
calibration, target perception, initialization uncertainty, appropriate vehicle
geometry and a domain-compliant validation trajectory remain outside this test.

\section{Reproduction and retained artifacts}
The local execution directory is
\path{simulation_output/carla_estimator_smoke_20261002/}. It retains
\path{captureSmoke.py}, \path{replaySmoke.m}, raw \path{capture/frames.jsonl},
logs and the original generated summaries. Execution entry points were:
\begin{quote}\small
\textbf{Regression:} MATLAB batch invocation of
\path{runtests('tests/onlineNrmmTrackingRuntimeTest.m')}, CSV export and
\texttt{assertSuccess(results)}.\par
\textbf{Capture:} \path{/home/zai/carla-venv/bin/python} invoked
\path{simulation_output/carla_estimator_smoke_20261002/captureSmoke.py}
with port 2011 and owned server PID 2478952.\par
\textbf{Replay:} \texttt{matlab -singleCompThread -batch} added
\path{simulation_output/carla_estimator_smoke_20261002/} to the MATLAB
path and invoked \texttt{replaySmoke}.
\end{quote}
The owned server PID is run-specific and must not be reused blindly.
Exact commands, capture metadata, cleanup verification, summary JSON,
input audit, source hashes, the 31-test table and both error traces are kept
under \path{report/CARLA_ESTIMATOR_SMOKE_20261002/}.
Identified export copies of the raw capture and execution adapters are kept
in the shared external archive at
\path{05_Project_B_Collision_Avoidance/2026-10-02_CARLA_Estimator_Smoke/}.
Raw recordings, legacy tools, generated kernels, unrelated working-tree
changes and dependencies are deliberately excluded from this project commit.
\end{document}
